\PassOptionsToPackage{table}{xcolor}
\documentclass[10pt,a4paper]{article}
\usepackage[T1]{fontenc}
\usepackage[utf8]{inputenc}
\usepackage[spanish,es-noshorthands]{babel}
\usepackage{tgpagella}
\usepackage{amsmath,amssymb,booktabs,tabularx,array,longtable,enumitem,multicol}
\usepackage{xcolor,tikz}
\usepackage[american]{circuitikz}
\usepackage[most]{tcolorbox}
\usepackage[margin=1.65cm,headheight=16pt,headsep=9pt]{geometry}
\usepackage{fancyhdr,lastpage,hyperref}
\usetikzlibrary{arrows.meta,positioning,calc}
\definecolor{ucbblue}{RGB}{0,51,102}
\definecolor{muted}{RGB}{241,245,249}
\definecolor{alert}{RGB}{140,57,39}
\hypersetup{colorlinks=true,linkcolor=ucbblue,urlcolor=ucbblue}
\newcolumntype{Y}{>{\raggedright\arraybackslash}X}
\newcolumntype{P}[1]{>{\raggedright\arraybackslash}p{#1}}
\newtcolorbox{note}[1][]{colback=muted,colframe=ucbblue!50,boxrule=.4pt,arc=1mm,#1}
\newtcolorbox{caution}[1][]{colback=red!3,colframe=alert,boxrule=.7pt,arc=1mm,#1}
\setlength{\parskip}{4pt}
\setlength{\parindent}{0pt}
\setlist{nosep,leftmargin=*}
\pagestyle{fancy}\fancyhf{}\lhead{\textbf{UCB Tarija} | IMT-221}\rhead{Guía experimental revisada}\cfoot{\thepage}\renewcommand{\headrulewidth}{0.4pt}
\newcommand{\status}{\textbf{Estado:} diseño didáctico propuesto; verificar la fuente y el fabricante del CI antes de energizar.}
\begin{document}

\begin{center}
{\Large\bfseries\textcolor{ucbblue}{Laboratorio 5 — Rectificación de precisión}}\\[2mm]
{\large Guía de montaje, medición y análisis: diodo 1N4148 frente a superdiodo LM358N}\\
\small IMT-221 · Circuitos Electrónicos II · Documento ampliado desde W14-S02 y W15-S01
\end{center}
\begin{note}\status\ Esta guía define un \emph{montaje de referencia} para facilitar la práctica. Las medidas, el fabricante y la alimentación de cada equipo se registran en banco; las predicciones indicadas no son resultados experimentales. Se mantiene el informe IEEE y la evidencia \textbf{MEDIDO/SIMULADO/PREDICHO}.\end{note}
\section{Objetivos y criterios de logro}
Al terminar, el equipo deberá: (1) construir un rectificador de media onda convencional y un superdiodo de media onda no inversor; (2) identificar cuándo el diodo y el lazo de realimentación conducen; (3) comparar las salidas para entradas pequeñas; (4) explicar saturación y recuperación; (5) entregar datos identificables para Lab 5 en su informe IEEE. La simulación no se registra como ejecución física. Ningún montaje de este documento ha sido declarado ensayado o validado en banco.
\section{Materiales y condiciones de prueba}
\begin{tabularx}{\linewidth}{@{}P{3.2cm}Y@{}}\toprule
Elemento por equipo & Uso previsto \\ \midrule
LM358N DIP-8 (1 de los 2 disponibles) & Amplificador A del superdiodo; segundo canal sin usar, tratar de acuerdo con datasheet/ficha de laboratorio.\\
1N4148 (2 de los 4) & Uno por rectificador.\\
Resistencias de $10\,\mathrm{k}\Omega$ (2) & $R_L$ de cada circuito; comprobar su valor real.\\
Protoboard, puentes; 100 nF (desacoplo si la configuración lo requiere) & Montaje y desacoplo en la proximidad del CI, bajo revisión docente.\\
Fuente de laboratorio regulada y limitada en corriente, generador, osciloscopio, multímetro & A proporcionar o confirmar por laboratorio. \\ \bottomrule\end{tabularx}
\begin{caution}\textbf{No conectar todavía.} Se propone una alimentación \emph{bipolar regulada} de $V_{S+}=+5\,\mathrm{V}$ y $V_{S-}=-5\,\mathrm{V}$ respecto a GND \textbf{únicamente como condición nominal del ensayo de referencia}, sujeta a aprobación. Si el banco no dispone de esta fuente o el fabricante exacto del LM358N no ha sido comprobado, ejecutar primero la ruta de simulación y \emph{no} aplicar señales negativas con fuente simple referida a tierra. Los límites de modo común y excursión de salida no son rail-to-rail y deben verificarse.\end{caution}
\section{Identificación del LM358N DIP-8 y referencias}
Vista superior del encapsulado, con muesca orientada hacia arriba (confirmar con ficha del componente recibido):
\begin{center}\begin{tabular}{@{}cl|cl@{}}\toprule 1 & Salida A & 8 & $V_{S+}$ \\2 & Entrada inversora A ($-$) & 7 & Salida B \\3 & Entrada no inversora A ($+$) & 6 & Entrada inversora B \\4 & $V_{S-}$ & 5 & Entrada no inversora B \\ \bottomrule\end{tabular}\end{center}
\textbf{Referencia de medida:} GND significa el punto común de la fuente y del generador. Confirmar cómo se conecta la pinza de masa del osciloscopio (normalmente unida a tierra de protección): no colocarla arbitrariamente en $-5\,\mathrm{V}$ ni sobre un nodo de señal. Identificar el fabricante/código completo del chip antes del encendido.
\section{Esquemas definitivos de referencia y conexiones}
\subsection*{Circuito A — media onda convencional}
\begin{center}\begin{circuitikz}[scale=.95,transform shape]
\draw (0,0) node[ground]{} to[sV,l_={$v_i$}] (0,2.2) to[D,l=$D_1$] (3.3,2.2) to[short,-o] (4.1,2.2) node[right]{$v_{o,D}$};
\draw (3.3,2.2) to[R,l={$R_L=10\,\mathrm{k}\Omega$}] (3.3,0) node[ground]{};
\end{circuitikz}\end{center}
\textbf{Conexiones:} positivo de generador al ánodo de $D_1$; cátodo (lado de la banda) al nodo $v_{o,D}$; $R_L$ desde $v_{o,D}$ a GND; negativo/referencia del generador a GND. Canal CH1 del osciloscopio a $v_i$, CH2 a $v_{o,D}$, ambas masas a GND verificado.
\subsection*{Circuito B — media onda de precisión no inversora}
\begin{center}\begin{circuitikz}[scale=.95,transform shape]
\node[op amp] (op) at (0,0) {};
\draw (op.+) to[short,-o] ++(-1.0,0) node[left]{$v_i$ (pin 3)};
\draw (op.out) to[short,*-] ++(0.6,0) coordinate(va) node[above]{$v_a$ (pin 1)} to[D,l=$D_2$] ++(2.2,0) coordinate(vo) to[short,*-o] ++(.8,0) node[right]{$v_{o,P}$};
\draw (vo) to[R,l={$10\,\mathrm{k}\Omega$}] ++(0,-2) node[ground]{};
\draw (op.-) -- ++(0,2.0) -| (vo);
\draw (op.up) -- ++(0,.35) node[above]{$V_{S+}$ (pin 8)};
\draw (op.down) -- ++(0,-.35) node[below]{$V_{S-}$ (pin 4)};
\end{circuitikz}\end{center}
\textbf{Conexiones:} pin 3 recibe $v_i$; pin 1 es $v_a$ y conecta al ánodo de $D_2$; el cátodo es $v_{o,P}$; conectar pin 2 directamente a $v_{o,P}$; $R_L=10\,\mathrm{k}\Omega$ desde $v_{o,P}$ a GND; pin 8 a $V_{S+}$, pin 4 a $V_{S-}$. No unir el pin 2 con GND. Si se utiliza solo el canal A, conectar el canal B como seguidor en reposo (pin 5 a GND, pin 6 al pin 7), previa comprobación del fabricante; no dejar entradas flotantes. CH1 mide entrada y CH2 la salida respecto a GND. Para medir $v_a$, mover una sonda de punta tras desconectar alimentación y verificar que la referencia no cambió.
\section{Predicción por estados}
El rectificador convencional con modelo de umbral predice $v_{o,D}\approx\max(v_i-V_\gamma,0)$. $V_\gamma$ es una aproximación dependiente de corriente, temperatura y diodo; no existe un valor fijo universal.

Para el superdiodo: en el semiciclo positivo y mientras el op-amp opera linealmente, $D_2$ conduce, se cierra el lazo y $v_{o,P}\approx v_i$, aunque la salida interna requiere $v_a\approx v_{o,P}+v_D$. En el negativo, $D_2$ bloquea, $v_{o,P}\approx0$ a través de $R_L$, el lazo se abre y la salida interna puede acercarse al límite negativo. No usar $v_+=v_-$ en el estado abierto. La recuperación de saturación y el cruce por cero pueden distorsionar la señal.
\begin{center}\begin{tabularx}{\linewidth}{@{}P{1.7cm}P{2cm}P{2.5cm}Y@{}}\toprule
Entrada & Diodo $D_2$ & Realimentación & Salida esperada y precaución \\ \midrule
$v_i>0$ & Conduce, si hay margen & Cerrada & $v_{o,P}\simeq v_i$ en región lineal; verificar $v_a$.\\
$v_i<0$ & Bloquea & Abierta & $v_{o,P}\simeq0$; salida interna puede saturar.\\
Cruce $v_i\approx0$ & Cambia de estado & Se abre/cierra & Buscar retardo/distorsión, offset y recuperación. \\ \bottomrule\end{tabularx}\end{center}
\section{Montaje progresivo: lista de autorización}
\begin{enumerate}[label=\textbf{Paso \arabic*.}]
\item \textbf{Sin alimentación:} fotografiar componentes, identificar bandas de diodos, muesca y pines del CI; comprobar con multímetro $R_L$ y ausencia de cortocircuito entre rieles.
\item Montar primero el convencional. Comprobar continuidad y orientación. Conectar generador y osciloscopio a GND común.
\item Montar el superdiodo \emph{sin energizar}. Revisar pin 2 unido exclusivamente al cátodo/salida, pin 3 al generador, pin 1 al ánodo; alimentar 8/4 según rieles aprobados.
\item Antes de encender: documentar tensión/corriente limitada, rango de modo común del datasheet exacto, rango previsto de $v_a$ y condiciones de seguridad del instrumento. Solicitar visto bueno del docente.
\item Encender fuente con generador inicialmente apagado y observar consumo/salida DC. Ante calentamiento o lectura incoherente, apagar y revisar, no manipular cables con el circuito energizado.
\item Aplicar senoidal de $100\,\mathrm{Hz}$, offset $0\,\mathrm{V}$ y amplitudes pico \emph{propuestas} $0.10$, $0.30$ y $0.80\,\mathrm{V}$, \textbf{solo tras verificar compatibilidad del banco}. Si se especifica Vpp en el generador, convertir $V_{pp}=2V_p$ y verificar la amplitud real sobre la carga de salida del generador.
\item Registrar simultáneamente $v_i$ y $v_o$ para ambas configuraciones; repetir una entrada a otra frecuencia opcional (p.ej. $1\,\mathrm{kHz}$) y justificar variaciones. Señales y capturas deben mostrar escala, acoplamiento, cero y origen MEDIDO/SIMULADO.
\end{enumerate}
\section{Plan mínimo de medidas}
\begin{tabularx}{\linewidth}{@{}P{1.4cm}P{1.5cm}P{1.7cm}P{1.6cm}P{1.6cm}Y@{}}\toprule
$V_{i,p}$ & $f$ & $V_{o,D,p}$ & $V_{o,P,p}$ & Fuente dato & Forma de onda, cruce por cero, archivo \\ \midrule
0.10 V & 100 Hz & & & & \\ \midrule
0.30 V & 100 Hz & & & & \\ \midrule
0.80 V & 100 Hz & & & & \\ \midrule
Elegido & Elegido & & & & \\ \bottomrule\end{tabularx}
\textbf{Datos:} colocar unidades en todas las casillas; si no hay dato, escribir \texttt{NO OBTENIDO} y motivo. Estas amplitudes son candidatos de prueba, no resultados ni condiciones impuestas cuando no sean seguras.
\section{Análisis y pregunta individual}
Definir $v_{\mathrm{ref}}(t)=\max(v_i(t),0)$, $e(t)=v_o(t)-v_{\mathrm{ref}}(t)$ y $E_{\mathrm{MAE}}=(1/T)\int_{t_0}^{t_0+T}|e(t)|\,dt$ [V]. Para muestras uniformes, $E_{\mathrm{MAE}}\approx N^{-1}\sum_{k=1}^N |e_k|$. Comparar \emph{la misma ventana, el mismo estímulo y la misma escala} en ambos circuitos. Distinguir error por umbral de error por saturación/offset/dinámica.
\begin{enumerate}
\item Explicar por qué el superdiodo aproxima la media onda ideal solo durante la conducción y dentro de límites del CI.
\item Identificar un punto de la señal donde la hipótesis de lazo cerrado deja de cumplirse.
\item Proponer una prueba para comprobar si la distorsión cerca de cero proviene del estado de salida interna o del osciloscopio.
\end{enumerate}
\section{Entregables IEEE y condiciones de aceptación}
Esquema real anotado, BOM efectivamente usada, foto del montaje, ficha de fabricante, protocolo de alimentación, tabla de señales con archivos, cálculo de error, discusión de saturación/cruce por cero y conclusión restringida a datos. Se adjunta la hoja W15-S01 de registro como apéndice opcional. \textbf{La práctica física se registra únicamente si fue ejecutada y supervisada}; la ruta simulada debe quedar identificada como tal.
\section*{Referencias y control de versión}
Basado en los archivos suministrados \texttt{w14\_s02\_guia\_estudiante\_lab\_5.tex} y \texttt{w15\_s01\_hoja\_registro\_lab\_5.tex}; materiales según BOM comunicada. Límites/pinout: verificar hoja de datos del \emph{LM358N del fabricante exacto}; como referencia técnica: TI, \emph{LMx58-N Low-Power Dual Operational Amplifiers}, \url{https://www.ti.com/lit/ds/symlink/lm358-n.pdf}. Las cifras de ensayos son propuestas, no mediciones realizadas.
\end{document}
